%% Chapter 6 — Section 6.2: Metropolis Adjusted Langevin Algorithm
\section{Metropolis Adjusted Langevin Algorithm}
\label{sec:6.2}

In the previous section, we have shown that the limit distribution of the Markov chain
$\{X^{(n)}\}_{n=1}^N$ defined by ULA does not agree with the target, even in a simple
Gaussian setting. In other words, the target $f \propto \exp(-V)$ is not an invariant
distribution for the Markov kernel
\[
  q_{\mathrm{ULA}}(x,\,\cdot\,) = \Normal\!\bigl(x - \epsilon\nabla V(x),\, 2\epsilon I\bigr)
\]
that defines the distribution of $X^{(n+1)}$ given that $X^{(n)} = x$. However, the proof of
Theorem~\ref{thm:6.2} suggests that, for small step-size $\epsilon$, the limit distribution
$f_\epsilon$ of ULA is of order $\epsilon$ apart from $f$, so that the kernel $q_{\mathrm{ULA}}$
\emph{approximately} preserves the target. This motivates the idea of using $q_{\mathrm{ULA}}$
as a proposal Markov kernel within the Metropolis Hastings algorithm studied in Chapter~4. The
key advantage of this choice of proposal is that since $q_{\mathrm{ULA}}$ approximately
preserves the target, the accept/reject mechanism in the Metropolis Hastings framework only
needs to account for a small correction to ensure convergence to the target. The resulting
method, known as the Metropolis Adjusted Langevin Algorithm (MALA), is summarized below.

\begin{algorithm}[H]
  \caption{Metropolis Adjusted Langevin Algorithm (MALA)}
  \label{alg:6.2}
  \begin{algorithmic}[1]
    \Require Target $f(x) \propto \exp(-V(x))$, initial distribution $\pi_0$,
             step-size $\epsilon > 0$, sample size $N$.
    \State Define the Markov kernel
    \[
      q_{\mathrm{ULA}}(x,\,\cdot\,) := \Normal\!\bigl(x - \epsilon\nabla V(x),\, 2\epsilon I\bigr).
    \]
    \State Run Metropolis Hastings (Algorithm~4.1) with inputs $f$, $\pi_0$, $q_{\mathrm{ULA}}$, $N$.
    \Ensure Sample $\{X^{(n)}\}_{n=1}^N$.
  \end{algorithmic}
\end{algorithm}

\begin{remark}[Pros and Cons]
MALA removes the bias of ULA at the price of introducing a Metropolis Hastings accept/reject
step. MALA shares many of the pros and cons of ULA. They are both well suited to sample
log-concave targets but can struggle to explore multi-modal distributions. Like ULA, MALA
requires evaluating $\nabla V$, which can be expensive. Compared to RWMH proposals from
Chapter~4, MALA can scale better to high-dimensional problems, particularly so for log-concave
targets. However, tuning the step-size parameter $\epsilon$ is more delicate than for RWMH;
existing theory shows that the step-size should be smaller in the transient than in the
stationary phase of the chain.
\end{remark}

\paragraph{Choice of Step-Size}
As we saw in Section~\ref{sec:6.1}, the limit distribution of ULA iterates depends on the
choice of step-size. The ability of ULA to produce samples that are approximately distributed
according to the target hinges on convexity of $V$ and on choosing a small step-size. On the
other hand, MALA corrects the bias of ULA via a Metropolis Hastings accept/reject mechanism.
For MALA, using a larger step-size leads to more rejections but faster exploration of the state
space. As for ULA, time-varying step-sizes can be used.

\paragraph{Scaling}
Classical theoretical results for MALA developed in the same large-$d$ scaling setting that we
considered for RWMH in Chapter~4 revealed two key insights.
\begin{enumerate}
  \item The step-size $\epsilon$ should be taken to decrease with $d$ at rate $d^{-1/3}$ in
    order to ensure that the acceptance rate remains bounded away from 0 and 1. Consequently,
    the number of iterations needed to get close to equilibrium is $\mathcal{O}(d^{1/3})$, a
    significant improvement over $\mathcal{O}(d)$ for RWMH. Though this fluctuates depending
    on the problem, many problems will have a cost per iteration of $\mathcal{O}(d)$ in both
    cases, so the overall complexity will be $\mathcal{O}(d^{4/3})$ compared to
    $\mathcal{O}(d^2)$.

  \item The optimal acceptance rate for MALA is $\approx 0.574$, compared to the considerably
    lower $\approx 0.234$ for RWMH.
\end{enumerate}

More recent results show enhanced scalability of MALA under various convexity assumptions on
the potential $V$, see Section~\ref{sec:6.4}. A heuristic reason for the improvement of MALA
over RWMH is that MALA, unlike RWMH, incorporates the target distribution in the proposal
kernel through the gradient $\nabla V$.
